\chapter{Proposal Generation and Mathematical Discovery}
\label{ch:06}

The second part of the book concerns instruments, and the first instrument to be examined is the one that has attracted the most attention and the least precision, namely the generation of proposals. A proposal in the relevant sense is a candidate answer to a well-posed problem in the sense defined in Chapter 4: a conjecture to be proved, a molecule to be synthesized, a design to be built, a policy to be adopted. The central fact about the generation of proposals by a computational system is that it is cheap relative to the evaluation of proposals, and that the relation between the two costs is likely to become more extreme as generation improves. A system that can produce a thousand plausible designs for a heat exchanger in the time that a team of engineers requires to produce one has not thereby produced a thousand heat exchangers, and it has produced no more than a set of hypotheses, each of which must still be tested against a standard that does not care how the hypothesis was obtained. The book calls the failure to keep this fact in view the error of premature adoption, and it regards it as the most likely practical mistake in the use of capable systems, more likely than the dramatic failures that dominate public discussion, because it is the natural consequence of an abundant supply of fluent and confident output.

The reasons for the persistent attractiveness of this error are not mysterious. Generated proposals arrive in the form that a competent proposal would take, with appropriate vocabulary, plausible reasoning, and apparent attention to objections, and the surface properties that human reviewers use as signals of quality are precisely the properties that a system trained on the products of human experts will reproduce. The interview that prompted this book draws attention to a related point in the context of language models, which are described as producing sequences of text that are probable given their training data and not as reporting facts, so that the fluency of the output is no evidence of its correctness. The point extends beyond language. A model that proposes molecules by interpolating among known compounds will produce molecules that look like known compounds, and a model that proposes proofs by imitating the style of published proofs will produce arguments that look like proofs. In each case the resemblance to the genuine article is a reason to subject the product to a test that depends on nothing about its appearance.

Mathematics provides the clearest case of how generation and verification can be separated, and it deserves detailed treatment for that reason. A mathematical proof is a finite object that can be checked, step by step, against the rules of a formal system. The checking is mechanical, it can be performed by a program of small size whose correctness can be examined independently, and it does not require any understanding of how the proof was found. This has the consequence that a mathematical result produced by a system of unknown reliability can nevertheless be accepted with the confidence appropriate to the checker, which is high. The history of the subject contains precedents. The proof of the four colour theorem by Appel and Haken relied on extensive computation that could not be verified by hand, and its acceptance was accompanied by a debate about what it means to know something that cannot be surveyed by a person; the debate was in a sense resolved by the later production of formally verified versions. The proof of the Kepler conjecture on sphere packing by Hales, which depended on large computations, was accompanied by a project to produce a formal proof, completed in 2014, whose purpose was exactly to replace the judgment of referees who had said they were unable to certify the original with a mechanical certificate. These precedents show that the verification of results of great length is practicable when the results are formalized, and they suggest that the contribution of capable systems to mathematics is likely to be assessed by the number of formally verified results they produce and not by the persuasiveness of their informal arguments.

It would nevertheless be a mistake to infer that the contribution is merely mechanical. Mathematics is not only the proof of stated theorems. It includes the formulation of conjectures, the invention of definitions, and the recognition that a certain structure is the right one for a class of problems, and these activities are judged by their fruitfulness, which is a human assessment that cannot be reduced to checking. A system that proposes conjectures from patterns in numerical data can contribute to the first of these, and there are documented cases in which such patterns have directed human mathematicians toward relationships they then proved. The role of the system in such cases is that of an instrument that makes certain regularities visible, and the human contribution lies in the interpretation and in the proof. The division of labor is stable and productive, and it is the pattern that the book proposes for other fields: the system generates, the verifier decides, and the human community determines what is worth asking.

For fields without a formal verifier the analogue of the proof assistant must be constructed, and the construction is an ordinary scientific and institutional task. A proposal for a new battery electrolyte is verified by synthesizing it and measuring its conductivity, stability, and cycle life, preferably in a laboratory that did not generate the proposal and that has no interest in its success. A proposal for a public health intervention is verified by a trial whose design was fixed before its results were known. A proposal for a bridge is verified by structural analysis, by testing of components, and by the independent review required by building codes. In each case the structure of the argument is the same, and the practical question is how to ensure that the verification is independent enough of the generation to be informative. The condition developed in Chapter 10 is that the errors of the verifier should not be correlated with the errors of the generator, since a verifier that shares the generator's blind spots confirms the very mistakes it is supposed to catch. This condition is violated in a subtle and important way when the verifier is another instance of the same kind of system, trained on the same data, and the book therefore treats the use of a model to grade a model as a screening step of limited weight and not as a substitute for external test.

A related consideration is the manner in which proposals are presented. A system that supplies a proposal together with a statement of confidence, an account of the evidence it considers relevant, and an explicit list of the assumptions on which the proposal depends is more useful than one that supplies the proposal alone, provided that the statements are themselves calibrated. Calibration is an empirical property that can be measured: among the proposals to which a system assigns probability of success $p$, a fraction approximately equal to $p$ should in fact succeed. A system that is badly calibrated, whether overconfident or underconfident, supplies information whose value must be discounted, and the calibration of the system on past problems in the same domain should be reported to those who decide. Proposals should also be accompanied by the record of the search that produced them, so that selection effects can be assessed: the best of a thousand candidates is a different object from a candidate that was the only one examined, and its apparent merit in the training or test data must be corrected for the number of alternatives considered.

The practical recommendation that follows is simple to state and demanding to implement. Each domain in which generation is to be accelerated should be accompanied by a plan for accelerating verification to a comparable degree, and the plan should specify the independent party that will perform verification, the tests it will use, the standard it will apply, and the sources of funding that insulate it from the success of the proposals it judges. Where the plan cannot be made, the rate of adoption should be set by the rate of verification and not by the rate of generation. This is not a counsel of timidity. It is the observation that the benefit of faster generation is realized only to the extent that its output is accepted by a process that can be trusted, and that a process that cannot be trusted converts abundance of proposals into abundance of risk.

\section*{Formalization}

The effect of verification on the reliability of accepted proposals can be computed. Suppose a generator produces proposals of which a fraction $q\in(0,1)$ satisfy the acceptance predicate $\varphi$ of a well-posed problem. A screening test is applied with sensitivity $s=\Pr(\text{pass}\mid\varphi=1)$ and false-positive rate $f=\Pr(\text{pass}\mid\varphi=0)$, with $0<f<s\le1$. Define the \emph{precision} of the test as $\pi=\Pr(\varphi=1\mid\text{pass})$.

\begin{proposition}
The precision after one test is $\pi=qs/(qs+(1-q)f)$. If $k$ tests are applied whose outcomes are conditionally independent given $\varphi$, and a proposal is accepted only if it passes all of them, the posterior odds are $\pi_k/(1-\pi_k)=\dfrac{q}{1-q}\left(\dfrac{s}{f}\right)^k$ when each test has sensitivity $s$ and false-positive rate $f$. Hence $\pi_k\to1$ as $k\to\infty$ at a geometric rate if and only if $s>f$, and the number of independent tests required for precision at least $\pi^\ast$ is
\[
k\;\ge\;\frac{\ln\bigl(\pi^\ast/(1-\pi^\ast)\bigr)-\ln\bigl(q/(1-q)\bigr)}{\ln(s/f)} .
\]
\end{proposition}

\begin{proof}
The first claim is Bayes' rule. For the second, conditional independence gives $\Pr(\text{all pass}\mid\varphi=1)=s^k$ and $\Pr(\text{all pass}\mid\varphi=0)=f^k$, and the posterior odds are the prior odds multiplied by the likelihood ratio $(s/f)^k$. The remaining statements follow by solving for $k$.
\end{proof}

The numbers show why the base rate $q$ governs the design of verification. For a hard domain in which only one proposal in a thousand is correct, $q=0.001$, a single test with $s=0.9$ and $f=0.01$, which would be regarded as a good screen, yields $\pi=0.0009/(0.0009+0.00999)\approx0.083$: more than ninety per cent of accepted proposals are wrong. Two independent such tests give odds of $(0.001/0.999)\cdot90^2\approx8.11$ and precision $0.890$, and three give precision $0.9986$. The proposition also exposes the cost of dependence among tests. If the tests are positively correlated given $\varphi=0$, the factor $f^k$ in the denominator is replaced by a larger quantity and the gain from each additional test is reduced. In the limit in which the tests are perfectly correlated the gain from a second test is zero, which is the formal statement of why a verifier that shares the generator's errors adds nothing.

The same structure supports a statement about calibration. Let $\hat p$ be the confidence reported by a generator for a proposal and let $\varphi\in\{0,1\}$ be the outcome. The generator is \emph{calibrated} if $\Pr(\varphi=1\mid\hat p)=\hat p$ almost surely. The \emph{calibration error} $\mathrm{CE}=\E\,|\Pr(\varphi=1\mid\hat p)-\hat p|$ can be estimated by binning on a validation set of proposals with known outcomes, and a reported confidence is to be used in decisions only after correction by the estimated map $\hat p\mapsto\Pr(\varphi=1\mid\hat p)$. Because that map is estimated from past problems, its validity is limited to problems of the same kind, which is the formal content of the instruction that the calibration of a system be reported by domain.
